\section{Block ground states of arbitrary positive local interactions}

The block-separation argument applies to any positive local interaction,
including interactions whose kernels are not the canonical local MPS spaces.
We follow the cyclic-cut argument of
\cite[Theorem~11, proof at lines~1418--1422]{PerezGarcia2007Matrix}.

\begin{lemma}[The periodic state in cyclic active-block coordinates]
    \label{lem:cyclic_mpv_active_boundary}
    \lean{MPSTensor.mpv_cyclicActiveBlockConfigEquiv_symm,
        MPSTensor.cyclicActiveBlock_mpvState_apply}
    \leanok
    \uses{def:ground_space_map, def:mpv}
    Let $A$ be a tensor, $L\le N$, and fix a cyclic cut $i$ of an $N$-site
    ring. If $\sigma$ is the active word of length $L$ and $\tau$ is the
    complementary word of length $N-L$, then the periodic coefficient,
    in these cyclic coordinates, is
    \begin{align}
        V^{(N)}(A)(\sigma,\tau)
         & =\tr(A^\sigma A^\tau)=\Gamma_L^A(A^\tau)(\sigma).
        \label{eq:cyclic_mpv_active_boundary}
    \end{align}
    The same identity holds for the Hilbert-space coordinate isometry.
\end{lemma}
\begin{proof}\leanok
    Rotate the closed word to begin at $i$, use invariance of the trace under
    this rotation, and split the resulting word into its two consecutive parts.
\end{proof}

\begin{lemma}[A translated interaction in active-block coordinates]
    \label{lem:cyclic_active_local_interaction}
    \lean{MPSTensor.cyclicActiveBlock_periodicLocalInteractionES_apply}
    \leanok
    \uses{def:periodic-fixed-interaction, lem:cyclic_mpv_active_boundary}
    Let $L\le N$, let $h$ be linear on $L$ sites, and fix a cyclic cut $i$.
    Let $U_i$ be the coordinate isometry of
    Lemma~\ref{lem:cyclic_mpv_active_boundary}, which splits the active block
    at $i$ from its complement. Then
    \begin{align}
        U_i(h_i\psi) & =(h\otimes I_{N-L})\,U_i(\psi).
        \label{eq:cyclic_active_local_interaction}
    \end{align}
    Thus $h$ acts on each fixed complementary fiber separately.
\end{lemma}
\begin{proof}\leanok
    Conjugate the defining translated interaction by the same coordinate
    isometry and cancel the inverse isometry.
\end{proof}

\begin{lemma}[A local operator separates the weighted block summands]
    \label{lem:arbitrary_local_block_separation}
    \lean{MPSTensor.periodicLocalInteractionES_smul_mpvState_eq_zero_of_wordTupleSpanTop}
    \leanok
    \uses{lem:cyclic_mpv_active_boundary, lem:cyclic_active_local_interaction,
        def:periodic-fixed-interaction}
    Let $A_1,\ldots,A_b$ be tensors with arbitrary finite virtual dimensions,
    let $c_j\in\C$, and let $h$ be any linear operator on $L$ sites of an
    $N$-site ring, with $L\le N$. Assume that the complementary word tuples
    $(A_1^\tau,\ldots,A_b^\tau)$, $|\tau|=N-L$, span the product of the
    block matrix algebras. If the translate $h_i$ annihilates
    $\sum_j c_jV^{(N)}(A_j)$, then it annihilates every
    $c_jV^{(N)}(A_j)$. The coefficients may vanish, and positivity of $h$
    is not needed for this local assertion.
\end{lemma}
\begin{proof}\leanok
    \uses{lem:cyclic_mpv_active_boundary, lem:cyclic_active_local_interaction,
        lem:local_weighted_boundary_annihilation}
    In the cyclic coordinates at $i$, compose the weighted joint boundary map
    with $h$: for boundaries $X=(X_1,\ldots,X_b)$ in the product algebra put
    \begin{align}
        T(X) & =h\Bigl(\sum_jc_j\,\Gamma_L^{A_j}(X_j)\Bigr).
        \label{eq:arbitrary_local_block_separation_map}
    \end{align}
    By (\ref{eq:cyclic_mpv_active_boundary}) and
    (\ref{eq:cyclic_active_local_interaction}), the hypothesis
    $h_i\sum_jc_jV^{(N)}(A_j)=0$ says $T(A_1^\tau,\ldots,A_b^\tau)=0$ for
    every complementary word $\tau$. These tuples span the product algebra, so
    $T=0$. Evaluating on the boundary $\delta_{jk}X$ supported in block $j$
    gives
    \begin{align}
        c_j\,h\,\Gamma_L^{A_j}(X) & =0
        \label{eq:arbitrary_local_block_separation_single}
    \end{align}
    for every $X$, which at $X=A_j^\tau$ is the local annihilation equation
    $h_i\,c_jV^{(N)}(A_j)=0$ of that block.
\end{proof}

\begin{lemma}[Arbitrary positive interactions preserve all block ground states]
    \label{lem:arbitrary_positive_block_ground_states}
    \lean{MPSTensor.periodicInteractionHamiltonianES_mpvState_eq_zero_of_wordTupleSpanTop}
    \leanok
    \uses{lem:arbitrary_local_block_separation, def:periodic-fixed-interaction}
    Under the complementary-span hypothesis of
    Lemma~\ref{lem:arbitrary_local_block_separation}, suppose $h\ge0$ and
    $c_j\ne0$ for every $j$. For $H=\sum_i h_i$,
    \begin{align}
        H\sum_j c_jV^{(N)}(A_j)=0
        \quad\Longrightarrow\quad
        HV^{(N)}(A_j)=0\quad\text{for every }j.
        \label{eq:arbitrary_positive_block_ground_states}
    \end{align}
\end{lemma}
\begin{proof}\leanok
    \uses{lem:arbitrary_local_block_separation, thm:weighted_sum_kernel}
    Each $h_i$ is positive. A finite sum of positive operators has kernel equal
    to the intersection of their kernels. Apply
    Lemma~\ref{lem:arbitrary_local_block_separation} to each $h_i$, cancel
    $c_j\ne0$, and sum the resulting equations.
\end{proof}

\begin{lemma}[Independence of the periodic block vectors]
    \label{lem:periodic_block_linear_independence}
    \lean{MPSTensor.linearIndependent_mpvState_of_wordTupleSpanTop}
    \leanok
    Let every block have positive virtual dimension. If the length-$N$ word
    tuples span the product matrix algebra, then the vectors
    $V^{(N)}(A_j)$ are linearly independent.
\end{lemma}
\begin{proof}\leanok
    A relation with coefficients $c_j$ says that the trace pairing against
    every word tuple vanishes for the boundary $(c_j I)_j$. The span forces
    each $c_jI=0$, and the positive block dimensions give $c_j=0$.
\end{proof}

\begin{theorem}[Degeneracy for separated canonical blocks and arbitrary interactions]
    \label{thm:canonical_arbitrary_interaction_degeneracy}
    \lean{MPSTensor.PGVWC07CanonicalFormData.periodicInteractionHamiltonianES_blocks_eq_zero,
        MPSTensor.PGVWC07CanonicalFormData.blockCount_le_finrank_ker_periodicInteractionHamiltonianES}
    \leanok
    \discussion{7684}
    \uses{def:pgvwc07_ti_canonical_representation,
        def:l_blk_injective, def:blocks_not_gauge_phase_equiv,
        lem:pgvwc07_block_tuple_span,
        lem:arbitrary_positive_block_ground_states,
        lem:periodic_block_linear_independence}
    Let $A$ have $b\ge2$ canonical blocks $A_j$ in the sense of
    \cite[Theorem~4]{PerezGarcia2007Matrix}, with positive block weights.
    Suppose no two distinct blocks are equivalent under a gauge and a phase,
    and suppose their words of a common length $L_0>0$ span their respective
    matrix algebras. For any positive operator $h$ on $L$ sites, put
    $H=\sum_i\tau^i(h)$ on an $N$-site ring. If
    \begin{align}
        N & \ge3(b-1)(L_0+1)+L,
          & HV^{(N)}(A)         & =0,
        \label{eq:canonical_arbitrary_interaction_degeneracy_hypotheses}
    \end{align}
    then
    \begin{align}
        HV^{(N)}(A_j) & =0\quad(1\le j\le b),
                      & \dim\ker H            & \ge b.
        \label{eq:canonical_arbitrary_interaction_degeneracy}
    \end{align}
    \textbf{Scope restriction (separated canonical blocks):}
    the gauge-phase separation premise is stronger than the pairwise
    inequality of finite-ring component vectors in the standing assumptions
    of \cite[Theorem~11]{PerezGarcia2007Matrix}; see
    \cite{gap:pgvwc07_direct_sum_input}.
\end{theorem}
\begin{proof}\leanok
    \uses{lem:pgvwc07_block_tuple_span,
        lem:arbitrary_positive_block_ground_states,
        lem:periodic_block_linear_independence}
    Canonical normalization, its positive dual fixed points and common
    injectivity derive the simultaneous spans at lengths $N-L$ and $N$.
    The original vector is the sum of the block vectors with nonzero
    coefficients $\lambda_j^N$. Apply
    Lemma~\ref{lem:arbitrary_positive_block_ground_states}; then
    Lemma~\ref{lem:periodic_block_linear_independence} gives $b$ independent
    vectors in $\ker H$.
\end{proof}
